\scanpage{51}
\mspage{50}

\uline{Récapitulation}

Les axiomes algébriques essentiels, en termes de la donnée $(\unsure{\mathcal F|\mathscr F|F},\leq,\ll)$, sont récapitulés p. 23. Diagrammes d'implication, concernant les axiomes d'\deleted{\unsure{in|en}}compatibilité-disjonction, voir p. 22 ou p. 45, 46 (avec une \deleted{$\unsure{\Lambda|\wedge}$} partie $\Lambda\subset\unsure{M|\mathcal M}$ dense à la clef, qu'on \unsure{voudra} \unsure{tjrs|tjs} dense -- notion d'atelier $\Lambda$-spacieux ou $\Lambda$-polyédral? \ldots). Sous les noms ``spacieux'' ou ``polyédral'', j'implique maintenant l'axiome $\mathrm{At}\,C(\unsure{\mathrm{spac}|\mathrm{spa}})$ resp. $\mathrm{At}\,\deleted{\unsure{p|po}}C(\mathrm{pol})$.

Dans le cas d'une structure $(M,\leq,\ll,\mathrel{>\mkern-13mu<})$, les données sont \deleted{\unsure{peu|plus}} plus compliquées (et encore, j'ai supprimé $\widetilde{\mathcal F}\inserted{\subset\unsure{P|\mathcal P}(M)}$ comme donnée explicite !), mais le syst. d'axiomes pertinent semble plus simple, \unsure{dans|deux|donc} \deleted{\unsure{certains}} \unsure{puis|pour} certains paquets d'axiomes regroupés.

Il y a lieu de considérer $(M,\leq,\ll,\mathrel{>\mkern-13mu<})$ comme une structure algébrique \unsure{indépendant|indépendante}, on pourrait l'appeler un ``pré-atelier'' (ou ``algèbre \unsure{de|des} multistrates''\deleted{\unsure{)|\}}}, ou ``algèbre de strates''). Elle peut correspondre à \uline{plusieurs} ateliers (réalisés comme des $\unsure{\overline{\mathcal F}|\widetilde{\mathcal F}}\subset P(M)$)

